\subsubsection*{ERCPMP}

The ERCPMP (Endoscopic Recognition of Colorectal Polyps Morphology and Pathology) dataset is a specialized resource focusing on the morphological and histopathological characterization of colorectal polyps. It consists of 796 annotated colonoscopy images and 21 corresponding video sequences, collected from 191 patients at a single clinical center between 2014 and 2019.

What distinguishes ERCPMP is its rich set of labels, which include not only anatomical features such as location, size, and circumference, but also standardized morphological descriptors (Paris classification, LST type, Kudo pit pattern, and JNET classification) and histological diagnoses (e.g., tubular, villous, serrated polyps; low- and high-grade dysplasia; carcinoma). All annotations were created and verified by trained gastroenterologists and pathologists.

The dataset is well-structured, with a separate metadata file linking each sample to its detailed clinical and pathological information. Although no benchmark models or evaluation metrics are provided, ERCPMP is intended to support the training of diagnostic and classification algorithms in both supervised and educational contexts.

By combining endoscopic morphology with definitive histopathological labels, ERCPMP enables the development of clinically relevant AI systems aimed at polyp risk stratification and treatment planning—tasks that go beyond simple detection and demand fine-grained classification.
